%! Functions and Function Notation — Unit 1, Lesson 1.1 — Homework
% THE LESSON'S GRADED PRACTICE and the last component in the packet — exactly TWO pages, started
% in the Homework Launch (the last 3 minutes of the period) and finished at home. Every context
% here is new: the notes used a rideshare fare, a vending machine, a garage tariff and an
% unlabeled graph; the AP Practice used a reservoir and a shipping tariff; the homework uses a
% phone plan, a delivery fee, a pool admission, a new graph, and a new formula.
%   Page 1 — Part A (evaluating, 1–3) and Part B (piecewise, 4–5).
%   Page 2 — Part C (a graph, 6–7) and Part D (building a rule, 8–10).
% The diagnostic item is 4(b): D(3) sits exactly on the boundary — the lesson's crux, met alone.
% Problem 10 asks for the lesson's headline in writing.
%
% PAGE LOCKSTEP with the other half of the pair: work blocks are byte-identical, prose answers
% are \writespace{H}{…} (fixed height both ways), and each \blank sits at the END of a line so
% its \ans cannot rewrap anything. The explicit \newpage fixes the page break in both files.
\documentclass[10pt]{article}
\usepackage{precalculus-article}
\usepackage{precalculus-boxes}
\newcommand{\TallMath}[1]{$\displaystyle #1\rule[-1.4em]{0pt}{3.2em}$}

\begin{document}

\pageheader{Unit 1, Lesson 1.1}{Homework}

\vspace{0.12in}

\boxguard[8]
\begin{remindbox}
\small
\textbf{This is your graded homework.} Start it in class; finish it on your own by the due
date on your cover sheet. Show your work: fill the slot with parentheses before you simplify,
and on any piecewise function decide \emph{which rule owns the input} before you compute.
\end{remindbox}

\vspace{0.12in}

\begin{headlinebox}{lilac}
\textbf{\color{plum}Part A --- Evaluating from a formula}
\end{headlinebox}

\begin{enumerate}[leftmargin=1.9em, itemsep=10pt]

\item A phone plan charges a monthly fee plus a charge for each gigabyte of data:
      $C(g) = 15 + 8g$, where $C(g)$ is the monthly cost in dollars for $g$ gigabytes.

      \textbf{(a)} Write in words what $C(3)$ means --- no computing yet.
      \writespace{1.1cm}{}

      \textbf{(b)} Find $C(3)$.

\begin{work}
  C(3) &= 15 + 8(3) \\
       &= 15 + 24 = 39
\end{work}

      $C(3) = $ \blank{2.2cm}

\item \begin{minipage}[t]{0.44\linewidth}
      Let $h(x) = 2x^2 - 5x$. Find $h(-3)$.

\begin{work}
  h(-3) &= 2(-3)^2 - 5(-3) \\
        &= 2(9) + 15 \\
        &= 33
\end{work}

      $h(-3) = $ \blank{2.2cm}
      \end{minipage}\hfill
      \begin{minipage}[t]{0.50\linewidth}
      \textbf{3.}\ Let $q(t) = -t^2 + 6$. Find $q(-2)$ and $q(0)$. \textit{Careful: square
      first, then apply the negative.}

\begin{work}
  q(-2) &= -(-2)^2 + 6 \\
        &= -4 + 6 = 2
\end{work}

      $q(-2) = $ \blank{1.8cm} \qquad $q(0) = $ \blank{1.8cm}
      \end{minipage}
\setcounter{enumi}{3}

\end{enumerate}

\vspace{0.14in}

\begin{headlinebox}{lilac}
\textbf{\color{plum}Part B --- Piecewise functions}
\end{headlinebox}

\begin{enumerate}[leftmargin=1.9em, itemsep=9pt, start=4]

\item A courier charges a flat fee for short trips and a per-mile rate beyond that. For a
      delivery of $m$ miles, the fee in dollars is
      \[ D(m) = \begin{cases} 4 & 0 < m \le 3 \\[2pt] 4 + 1.5(m - 3) & 3 < m \le 12 \end{cases} \]

      \textbf{(a)} $D(2) = $ \blank{2.0cm} \hfill
      \textbf{(b)} $D(3) = $ \blank{2.0cm} \hfill
      \textbf{(c)} $D(9) = $ \blank{2.0cm}

      \textbf{(d)} Which rule owns $m = 3$, and how do you know? Name the condition you checked.
      \writespace{1.2cm}{}

\item Let $v(x) = x^2 + 1$ for $x \le -2$, and $v(x) = 3 - x$ for $x > -2$.

      $v(-4) = $ \blank{1.8cm} \qquad $v(-2) = $ \blank{1.8cm} \qquad $v(0) = $ \blank{1.8cm}

\end{enumerate}

\newpage

\begin{headlinebox}{lilac}
\textbf{\color{plum}Part C --- Reading a graph}
\end{headlinebox}

\begin{minipage}[t]{0.47\linewidth}
\vspace{0pt}
\centering
\begin{tikzpicture}[x=0.62cm, y=0.5cm, >=stealth]
  \draw[linegray, very thin, step=1] (-4,-2) grid (6,6);
  \draw[->, thick] (-4,0) -- (6.5,0) node[below right] {\small $x$};
  \draw[->, thick] (0,-2) -- (0,6.5) node[above] {\small $g(x)$};
  \foreach \x in {-3,-2,-1,1,2,3,4,5} \node[below, font=\scriptsize] at (\x,0) {\x};
  \foreach \y in {-1,1,2,3,4,5} \node[left, font=\scriptsize] at (0,\y) {\y};
  \draw[very thick, plum] (-3,3) -- (-1,-1) -- (2,5) -- (5,2);
  \fill[plum] (-3,3) circle (2.5pt);
  \fill[plum] (5,2) circle (2.5pt);
\end{tikzpicture}
\end{minipage}\hfill
\begin{minipage}[t]{0.5\linewidth}
\vspace{0pt}
\begin{enumerate}[leftmargin=1.9em, itemsep=10pt, start=6]
\item Find each value from the graph.

      $g(2) = $ \blank{1.8cm}

      $g(-1) = $ \blank{1.8cm}

      $g(5) = $ \blank{1.8cm}

\item State the domain and range of $g$ as inequalities, then again with brackets.

      Domain: \blank{5.4cm}

      Range: \blank{5.4cm}
\end{enumerate}
\end{minipage}

\vspace{0.16in}

\begin{headlinebox}{lilac}
\textbf{\color{plum}Part D --- Building a rule, and the headline}
\end{headlinebox}

\begin{enumerate}[leftmargin=1.9em, itemsep=10pt, start=8]

\item A city pool charges \$5 admission for visitors aged 12 and under, and \$9 for visitors
      over 12, up to an age limit of 80. Write $P(a)$, the admission in dollars for a visitor
      aged $a$, as a piecewise function with both conditions.
      \writespace{2.0cm}{}

\item Let $r(x) = \dfrac{12}{x + 4}$.

      \textbf{(a)} Find $r(2)$. \quad $r(2) = $ \blank{2.2cm}

      \textbf{(b)} Which input is forbidden, and why? Write the domain of $r$ in words.
      \writespace{1.4cm}{}

\item A classmate evaluates a piecewise function by working out \emph{both} rules and then
      picking whichever answer looks nicer. Use the words \emph{input} and \emph{condition} to
      explain why exactly one rule applies to any input.
      \writespace{1.9cm}{}

\end{enumerate}

\vspace{0.1in}

\boxguard[8]
\begin{spiralbox}
\small
\textbf{Next lesson (1.2, Function Operations):} today one function ran at a time. Next we add,
subtract, multiply, and divide two functions to build a third. Look back at problem 1: if a
second plan charged $S(g) = 5 + 2g$ for streaming, what would the two plans cost \emph{together}
for $g$ gigabytes?
\end{spiralbox}

\end{document}
